\documentclass[10pt,letterpaper]{article}
\usepackage[top=0.85in,left=2.75in,footskip=0.75in]{geometry}

\usepackage{amsmath,amssymb}
\usepackage{changepage}
\usepackage[utf8x]{inputenc}
\usepackage{textcomp}
\usepackage[T1]{fontenc}
\usepackage{cite}
\usepackage{nameref,hyperref}
\usepackage[right]{lineno}
\usepackage{microtype}
\DisableLigatures[f]{encoding = *, family = * }
\usepackage{xcolor}
\usepackage{graphicx}
\usepackage{booktabs}
\usepackage{multirow}
\usepackage{sectsty}
\allsectionsfont{\sffamily\Large}
\usepackage{fancyhdr}
\pagestyle{fancy}
\fancyhf{}
\rfoot{\thepage}
\renewcommand{\headrulewidth}{0pt}
\linenumbers

\newcommand{\beginsupplement}{%
  \setcounter{table}{0}
  \renewcommand{\thetable}{S\arabic{table}}%
  \setcounter{figure}{0}
  \renewcommand{\thefigure}{S\arabic{figure}}%
}

%% ============================================================
%% TITLE
%% ============================================================

\title{Condition-dependent binding QTL switching reveals the transcription
factor concentration landscape as a gatekeeper of functional
genetic variation in maize}

\author{
Michael~Banf\textsuperscript{1},
Thomas~Hartwig\textsuperscript{2,3,4,*}
}

\date{}

\begin{document}

\vspace*{0.2in}
\begin{flushleft}
{\Large\bfseries\sffamily Condition-dependent binding QTL switching reveals the
transcription factor concentration landscape as a gatekeeper of functional
genetic variation in maize}
\bigskip

\noindent
Michael~Banf\textsuperscript{1} and Thomas~Hartwig\textsuperscript{2,3,4,*}

\bigskip

\noindent
\textbf{1}~Perelyn, Munich, Germany\\
\textbf{2}~Heinrich Heine University D\"usseldorf, Faculty of Mathematics and Natural Sciences, Institute for Molecular Physiology, D\"usseldorf, Germany\\
\textbf{3}~Independent Research Groups, Max Planck Institute for Plant Breeding Research, Cologne, Germany\\
\textbf{4}~Cluster of Excellence on Plant Sciences (CEPLAS), Heinrich Heine University, D\"usseldorf, Germany

\bigskip

\noindent
*~Corresponding author.

\noindent
E-mail: hartwit@hhu.de

\end{flushleft}

% ============================================================
\section*{Abstract}
% ============================================================

Binding quantitative trait loci (bQTL) capture sequence variants that alter
transcription factor (TF) occupancy, yet whether a given bQTL produces a
functional allele-specific expression (ASE) effect is not a fixed property of
the variant. Here, we exploit paired allele-specific expression data from
25~maize F1 hybrids under well-watered and drought stress conditions to
demonstrate that bQTL functionality is condition-dependent: only 31\% of
functional bQTL are shared between conditions, while 37\% are well-watered--specific
and 31\% are drought-specific (Jaccard index = 0.314). Using NAM founder
genotypes at 43,900 bQTL positions, we establish a direct causal chain from
variant genotype to ASE magnitude (814 FDR-significant associations, 5.2-fold
enrichment over expectation) and show that this relationship is
condition-gated. Functional bQTL are distinguished by proximity to the
transcription start site (constitutive: 228\,bp vs non-functional: 579\,bp,
$p = 3.9 \times 10^{-13}$) and by the expression level of their target genes,
but---critically---\emph{not} by which TF family's motif is disrupted. Under
well-watered conditions, TF family expression is strongly anti-correlated with
bQTL effect size ($\rho = -0.830$, $p = 0.0008$), consistent with a saturation
model where abundant TFs mask allelic differences. Drought abolishes this
organized relationship ($\rho = -0.067$, $p = 0.84$), and drought-specific
bQTL are enriched at NAC-bound ($\text{OR} = 1.30$) and TCP-bound
($\text{OR} = 1.51$) sites whose cognate TFs are drought-induced. We propose a
\emph{regulatory configuration} model in which the TF concentration landscape
acts as a condition-dependent gate that determines which sequence variants
produce detectable functional consequences---the same variant, different
impact.


% ============================================================
\section*{Introduction}
% ============================================================

Genetic variation at transcription factor binding sites is a major source of
phenotypic diversity in plants and animals \cite{engelhorn2025,peleke2024}.
In maize, a recent pan-cistrome study identified 147,942 binding quantitative
trait loci (bQTL)---sequence variants that alter TF occupancy in at least one
of 25~F1 hybrids relative to the B73 reference \cite{engelhorn2025}. These
bQTL collectively explained 71\% of heritable trait variation, establishing
cis-regulatory variation as a dominant driver of phenotype.

Yet a fundamental question remains: are all bQTL equally functional, or does
their impact on gene expression depend on cellular context? The standard model
treats cis-regulatory variants as having fixed effect sizes determined by
sequence alone---a variant either disrupts a binding motif or it does not. This
view, however, ignores the fact that TF binding is a mass-action equilibrium:
the occupancy of a site depends not only on its affinity but also on the
concentration of the cognate TF \cite{segal2008,levo2012}. A high-affinity
variant in a motif for a TF that is not expressed will have no functional
consequence; conversely, a moderate-affinity variant may become critical when
its TF is abundant.

This concentration-dependent logic predicts that bQTL should exhibit
\emph{condition-dependent switching}: the same variant may be functionally
consequential under one condition but silent under another, depending on the
TF expression landscape. Such switching has been documented at individual loci
in model organisms \cite{kumasaka2016,ban2015,moyerbrailean2016}, but a
systematic, genome-wide test has not been performed in plants.

Here, we exploit a unique feature of the maize pan-cistrome dataset: Engelhorn
et al.\ collected allele-specific expression (ASE) data for all 25~F1 hybrids
under both well-watered (WW) and drought stress (DS) conditions
\cite{engelhorn2025}. Although the original study focused on bQTL
identification under control conditions, the paired ASE data enables a direct
test of condition-dependent bQTL functionality. We combine this with NAM
founder genotype data \cite{grzybowski2023} to establish a causal chain from
variant genotype through TF motif disruption to allele-specific expression,
and show that this chain is gated by the TF concentration landscape.


% ============================================================
\section*{Results}
% ============================================================

% ------------------------------------------------------------
\subsection*{Variant genotype causally determines allele-specific expression}
% ------------------------------------------------------------

To test whether bQTL genotype directly causes ASE differences among F1
hybrids, we obtained NAM founder genotypes at 43,900 bQTL positions from
whole-genome sequencing data \cite{grzybowski2023}. For each bQTL within
2\,kb of a gene's transcription start site (TSS), we split the 24 non-B73
hybrids into two groups---those whose NAM parent carries the variant allele
and those with the reference allele---and compared their signed ASE values
(log$_2$ B73/NAM allelic ratio).

Under well-watered conditions, 814 of 8,040 testable bQTL--gene pairs were
significant at FDR $< 0.05$ (Mann-Whitney $U$ test), representing a 5.2-fold
enrichment over chance expectation (Fig.~\ref{fig:causal}A). The
enrichment is not an artifact of multiple testing: the overall $p$-value
distribution shows a strong left skew (Fig.~\ref{fig:causal}B), and Storey's
$\pi_1$ estimate indicates that 37.6\% of tested pairs harbor true
associations. The median effect size among significant pairs is
$|d| = 2.13$ (Cohen's $d$); these large values reflect the small per-group
sample sizes ($n = 3$--$7$ hybrids) inherent in the split-by-genotype design,
which inflate point estimates while the statistical significance remains
robust.

Importantly, the test controls for all trans-acting variation: all hybrids
share B73 as one parent, so differences in ASE between the genotype groups can
only arise from the cis-acting variant itself.

% ------------------------------------------------------------
\subsection*{TF motif identity is irrelevant under baseline conditions}
% ------------------------------------------------------------

Given that bQTL typically disrupt TF binding motifs, we asked whether the
identity of the disrupted TF family predicts the magnitude of the genotype--ASE
effect. Using 259 position weight matrices from PlantTFDB~5.0 \cite{tian2020},
we scanned the $\pm$28\,bp region around each bQTL for motif overlaps and
classified bQTL by the TF families whose motifs are disrupted.

The result was unequivocal: under well-watered conditions, all TF families
produce nearly identical genotype--ASE effect sizes. The coefficient of
variation across 18 families is only 4.6\%, with family-level mean
$|d|$ ranging from 0.75 to 0.89 (Fig.~\ref{fig:motif_irrelevance}A).
Neither the number of disrupted motifs ($r_s = 0.02$, $p = 0.85$) nor the
specific family ($F$-test $p = 0.71$) predicts effect magnitude.

This irrelevance of TF identity is striking. It implies that under baseline
conditions, the regulatory system is sufficiently buffered that the specific
transcription factor whose motif is disrupted does not matter---what matters is
simply \emph{that} a motif is disrupted at a transcriptionally active locus.

% ------------------------------------------------------------
\subsection*{Regulatory grammar of functional bQTL}
% ------------------------------------------------------------

If TF identity does not distinguish functional from non-functional bQTL, what
does? We compared 814 functional (FDR $< 0.05$) with 7,226 non-functional
bQTL--gene pairs across multiple genomic and regulatory features.

Two features emerged as significant discriminators. First, \textbf{distance to
TSS}: functional bQTL are significantly closer to the transcription start site
(mean 521\,bp vs 649\,bp, $p = 1.2 \times 10^{-6}$, Mann-Whitney $U$).
Constitutive bQTL---those functional under both conditions---are closest
(median 228\,bp), followed by drought-specific (334\,bp), WW-specific
(446\,bp), and non-functional (579\,bp). This gradient suggests that proximal
promoter variants are always consequential, while distal variants require
additional context.

Second, \textbf{target gene expression level}: genes harboring functional bQTL
have lower mean expression (mean 698\,CPM vs 1,000\,CPM for non-functional,
$p = 0.007$), driven by a tail of highly expressed non-functional targets.
This is consistent with a model where highly expressed genes have saturated
TF occupancy that masks allelic effects.

Notably, motif redundancy (number of TF family motifs in the promoter),
number of bQTL per promoter, and the presence of complete regulatory chains
(expressed TF $\to$ disrupted motif $\to$ expressed target) did \emph{not}
significantly discriminate functional from non-functional bQTL, further
supporting the irrelevance of specific TF identity under baseline conditions.

% ------------------------------------------------------------
\subsection*{Condition-dependent bQTL switching}
% ------------------------------------------------------------

To test whether bQTL functionality is condition-dependent, we repeated the
genotype--ASE test using drought-stress ASE data from the same 25 hybrids
(Engelhorn et al.\ Supplementary Table S13b \cite{engelhorn2025}).

Drought produced a comparable overall signal: 742 FDR-significant pairs
(9.1\%, 4.95-fold enrichment, $\pi_1 = 0.369$). However, the
\emph{identity} of functional bQTL changed dramatically
(Fig.~\ref{fig:switching}A).

Of the 1,184 unique bQTL--gene pairs that are functional in at least one
condition:
\begin{itemize}
  \item 372 (31.4\%) are \textbf{constitutive}---significant in both WW and DS
  \item 442 (37.3\%) are \textbf{WW-specific}---significant only under
        well-watered conditions
  \item 370 (31.2\%) are \textbf{drought-specific}---significant only under
        drought
\end{itemize}

The Jaccard index of 0.314 indicates that the two conditions share less than
one-third of their functional bQTL, despite testing the same variants in the
same genotypes. This is not a statistical power issue: the overall enrichment
and $\pi_1$ are similar between conditions, and the number of FDR hits is
comparable (814 WW vs 742 DS).

Critically, the same genotypes are tested under both conditions
($|\Delta n_\text{variant}| \approx 0.1$ between conditions), confirming
that the switching is not driven by changes in genotype composition but by
changes in the \emph{regulatory environment}.

% ------------------------------------------------------------
\subsection*{Switching is a threshold effect, not binary on/off}
% ------------------------------------------------------------

Examination of paired effect sizes reveals that switching bQTL do not
abruptly gain or lose function---they modulate continuously
(Fig.~\ref{fig:switching}B,C).

WW-specific bQTL have strong effects under WW ($|d| = 1.989$) that
\emph{decrease by 50\%} under drought ($|d| = 0.985$). Conversely,
drought-specific bQTL have modest effects under WW ($|d| = 1.052$) that
\emph{increase by 83\%} under drought ($|d| = 1.920$). Constitutive bQTL
maintain the largest effects in both conditions ($|d| \approx 2.3$),
well above the significance threshold regardless of condition.

The directionality of effects is highly conserved: 98.9\% of constitutive,
92.6\% of WW-specific, and 91.6\% of drought-specific bQTL show the same
sign of Cohen's $d$ in both conditions. The variant is doing the same thing
in both conditions---what changes is whether the magnitude crosses the
detection threshold.

This continuous modulation rules out a model in which drought activates
entirely new regulatory programs. Instead, the data support a \emph{dimmer
switch} model: each bQTL has an intrinsic effect that is amplified or
attenuated by the regulatory environment.

% ------------------------------------------------------------
\subsection*{Target gene drought response predicts bQTL switching}
% ------------------------------------------------------------

If bQTL switching is not driven by the variant itself, it may be driven by
changes in the target gene's transcriptional context. We tested this by
comparing the drought expression response (log$_2$FC DS/WW) of target genes
across bQTL categories (Fig.~\ref{fig:switching}D).

Drought-specific bQTL target genes that are significantly more
drought-responsive than WW-specific bQTL targets (median log$_2$FC
$= +0.413$ vs $+0.219$, $p = 4.1 \times 10^{-4}$, Mann-Whitney $U$).
Furthermore, 37.0\% of DS-only target genes are strongly drought-responsive
($|$log$_2$FC$| > 1$) compared to 29.8\% of WW-only targets.

Constitutive bQTL targets are intermediate in drought response (median
log$_2$FC $= +0.335$) but, strikingly, are the \emph{least} likely to be
strongly drought-responsive (24.3\% with $|$log$_2$FC$| > 1$). This
seemingly paradoxical result makes sense in the dimmer switch model:
constitutive bQTL are close to the TSS (228\,bp) and have the largest
intrinsic effects ($|d| \approx 2.3$), so they remain above threshold
regardless of how the gene's expression program changes.

At the gene level, the switching is nearly complete: of 1,039 genes with
functional bQTL in either condition, only 341 (33\%) have functional bQTL
in both, while 386 are WW-only and 312 are DS-only. Only 29 genes exhibit
mixed switching (different bQTL functional in different conditions within
the same promoter).

% ------------------------------------------------------------
\subsection*{The TF concentration landscape gates bQTL functionality}
% ------------------------------------------------------------

The preceding results show that bQTL switching correlates with target gene
expression changes, but the mechanistic driver must be the TFs that read the
disrupted motifs. We therefore tested whether TF expression levels predict
bQTL effect sizes.

Under well-watered conditions, TF family expression (mean CPM across hybrids)
is \emph{strongly anti-correlated} with the genotype--ASE effect size
(Spearman $\rho = -0.830$, $p = 0.0008$; Fig.~\ref{fig:saturation}A). This
is the signature of a \textbf{saturation model}: when a TF is abundant, both
the high-affinity (reference) and reduced-affinity (variant) alleles are
occupied, masking the allelic difference. When a TF is scarce, only the
high-affinity allele is occupied, and the variant allele shows reduced binding
and expression.

Under drought, this organized relationship \emph{collapses}
($\rho = -0.067$, $p = 0.84$; Fig.~\ref{fig:saturation}B). The TF
concentration landscape is reorganized by stress: NAC family TFs increase
5.25-fold (log$_2$FC), ERF 0.51-fold, HSF 0.50-fold, while G2-like decreases
0.87-fold (Fig.~\ref{fig:saturation}C). This reorganization
disrupts the baseline concentration--effect relationship, creating a new
landscape in which different bQTL cross the detection threshold.

Per-TF-family analysis reveals two distinct behaviors under drought:
\begin{enumerate}
  \item \textbf{Upregulated TFs lose bQTL effect} (bHLH, ERF, HSF,
    Trihelix)---consistent with the saturation model. As these TFs become
    more abundant, they saturate both alleles, reducing the allelic difference.
  \item \textbf{Developmental TFs gain bQTL effect} (LFY $\Delta|d| = +0.25$,
    YABBY $+0.17$)---likely because their targets become more sensitive as
    developmental programs are reprioritized under stress.
\end{enumerate}

NAC presents an interesting exception: despite being the most upregulated
family (+5.25 log$_2$FC), NAC-associated bQTL show a slight \emph{increase}
in effect ($\Delta|d| = +0.014$). NAC TFs may have unique binding dynamics,
or the massive induction may activate NAC targets so strongly that even
partial binding differences become detectable.

% ------------------------------------------------------------
\subsection*{TF family enrichment at switching bQTL}
% ------------------------------------------------------------

If the TF concentration landscape gates bQTL functionality, we predict that
drought-specific bQTL should be enriched for motifs of drought-induced TFs.
We tested this by computing odds ratios for each TF family's representation
among WW-specific vs drought-specific bQTL (Fig.~\ref{fig:saturation}D).

Drought-specific bQTL are enriched for NAC motifs (OR $= 1.30$) and TCP
motifs (OR $= 1.51$), both families that are upregulated under drought.
Conversely, MYB\_related motifs are depleted from drought-specific bQTL
(OR $= 0.51$), consistent with MYB being a constitutively expressed family.
GRAS shows the strongest drought enrichment (OR $= 1.47$), and bZIP---a
known stress-responsive family---is also enriched (OR $= 1.32$).

These enrichments confirm the regulatory configuration model: drought-induced
TFs create a new concentration landscape that activates a distinct set of
bQTL---specifically, those at motifs for the induced TFs.


% ============================================================
\section*{Discussion}
% ============================================================

% ------------------------------------------------------------
\subsection*{The regulatory configuration model}
% ------------------------------------------------------------

Our results establish a new framework for understanding the functional
consequences of cis-regulatory genetic variation. The standard model treats
bQTL as having fixed effect sizes determined by sequence alone. We show
instead that bQTL functionality is \emph{gated} by the TF concentration
landscape---the same variant can have a large, small, or undetectable effect
depending on the expression levels of its cognate TFs and the transcriptional
activity of its target gene.

This regulatory configuration model has three tiers:

\begin{enumerate}
  \item \textbf{Core promoter variants} ($< 250$\,bp from TSS) are
    constitutively functional because they affect basal transcriptional
    machinery that is always present.
  \item \textbf{Proximal regulatory variants} (250--500\,bp) are
    condition-dependent, functional only when their cognate TFs are
    sufficiently expressed to create detectable allelic differences.
  \item \textbf{Distal regulatory variants} ($> 500$\,bp) require both
    TF expression and additional chromatin context to be functional.
\end{enumerate}

The gating mechanism is TF binding saturation: at high TF concentrations,
both alleles are occupied and allelic differences are masked; at intermediate
concentrations, only the high-affinity allele is occupied, maximizing the
allelic effect; at very low concentrations, neither allele is occupied and
the variant is silent.

% ------------------------------------------------------------
\subsection*{Implications for missing heritability}
% ------------------------------------------------------------

The condition-dependent switching we observe has direct implications for the
``missing heritability'' problem in quantitative genetics
\cite{manolio2009}. If bQTL effect sizes are not fixed but depend on the
regulatory configuration, then:

\begin{itemize}
  \item GWAS conducted under a single condition will capture only the subset
    of functional variants active in that condition
  \item Effect size estimates from one environment may not transfer to another
  \item $G \times E$ interactions for cis-regulatory variants are not
    exceptions but the \emph{norm}---69\% of functional bQTL in our data are
    condition-specific
\end{itemize}

This does not increase the total number of causal variants but redistributes
their effects across conditions, potentially explaining why multi-environment
GWAS recover more heritable variation than single-environment studies
\cite{deng2022}.

% ------------------------------------------------------------
\subsection*{The saturation model and TF concentration}
% ------------------------------------------------------------

The strong anti-correlation between TF expression and bQTL effect size
($\rho = -0.830$) under well-watered conditions provides direct evidence for
a thermodynamic binding equilibrium at regulatory variants
\cite{segal2008,levo2012}. This observation predicts that:

\begin{itemize}
  \item Rare TFs (low expression) should produce larger allelic effects
    because their binding sites are not saturated
  \item Stress-induced TFs should initially increase and then decrease bQTL
    effects as their concentration rises through the binding isotherm
  \item The ``irrelevance of TF identity'' under baseline conditions reflects
    a state where most TFs are at concentrations that produce similar
    fractional occupancy differences between alleles
\end{itemize}

The collapse of the concentration--effect relationship under drought
($\rho = -0.067$) indicates that the system has been pushed out of its
organized baseline into a reconfigured state where the binding equilibria
are no longer described by a single concentration--effect axis. This is
consistent with the known complexity of drought stress responses, which
involve chromatin remodeling, post-translational modifications of TFs, and
interactions between stress and developmental signaling pathways
\cite{kim2017,deng2022}.

% ------------------------------------------------------------
\subsection*{Comparison with Engelhorn et al.\ (2025)}
% ------------------------------------------------------------

Our analysis builds directly on the bQTL and ASE datasets generated by
Engelhorn et al.\ \cite{engelhorn2025} but addresses fundamentally different
questions. The original study focused on bQTL identification and their
contribution to heritable trait variation under control conditions. Although
drought ASE data were collected and included as supplementary material
(Tables~S13a,b), no cross-condition comparison of bQTL functionality was
performed.

Our key novel contributions are: (1)~a causal test linking variant genotype
to ASE magnitude using NAM founder genotypes; (2)~the discovery that TF motif
identity is irrelevant under baseline conditions; (3)~the demonstration that
69\% of functional bQTL are condition-specific; (4)~the TF saturation model
explaining condition-dependent switching; and (5)~the regulatory
configuration framework unifying these observations.

% ------------------------------------------------------------
\subsection*{Limitations}
% ------------------------------------------------------------

Several limitations warrant discussion. First, our genotype data cover only
29.7\% of bQTL positions (43,900 of 147,942), limiting the power of the
genotype--ASE test. Second, we lack allele-specific information for the bQTL
themselves---we know which position varies but not which allele increases or
decreases binding. Third, the drought treatment is a single time point, and
bQTL switching dynamics across the stress response are unknown. Fourth, our
TF expression estimates are gene-level averages across tissues, not
cell-type-specific measurements.

Despite these limitations, the consistency of the results across multiple
analyses---genotype--ASE tests, TF expression correlations, motif enrichments,
and genomic architecture---supports the robustness of the regulatory
configuration model.


% ============================================================
\section*{Methods}
% ============================================================

% ------------------------------------------------------------
\subsection*{bQTL and expression data}
% ------------------------------------------------------------

We used the 147,942 bQTL positions from Engelhorn et al.\ \cite{engelhorn2025}
(Nature Genetics Supplementary Table~S6). Allele-specific expression data
under well-watered and drought conditions were obtained from Supplementary
Tables S13a and S13b, respectively, which provide per-hybrid allele-specific
read counts for $\sim$25,000 genes. For each gene and hybrid, we computed the
allele-specific expression ratio as log$_2$(B73 reads / NAM reads),
averaging across three biological replicates per condition. Gene-level
expression data (well-watered vs drought, 24,796 genes) were from
Supplementary Table~S7.

% ------------------------------------------------------------
\subsection*{NAM founder genotypes}
% ------------------------------------------------------------

Whole-genome sequencing VCF files for the 26 NAM founders were obtained from
Grzybowski et al.\ \cite{grzybowski2023} via the CyVerse Data Store. Per-chromosome
VCFs were queried at all bQTL positions using bcftools
\texttt{view -R}, yielding genotypes for 43,900 bQTL (29.7\% of total)
across 24 non-B73 founders. Phased genotypes (0|0, 0|1, 1|1) were converted
to binary variant-carrier status, and founder genotypes were mapped to F1
hybrid identifiers via the known crossing scheme (B73 $\times$ NAM founder).

% ------------------------------------------------------------
\subsection*{Genotype--ASE causal test}
% ------------------------------------------------------------

For each bQTL within 2\,kb of a gene's TSS, hybrids were split into two
groups based on whether their NAM parent carries the variant allele. We
required $\geq 3$ hybrids per group and $\geq 3$ non-missing ASE values
per group. We tested for differences in signed ASE (log$_2$ B73/NAM) using
the Mann-Whitney $U$ test (two-sided) and Welch's $t$-test, with effect
sizes quantified as Cohen's $d$. Multiple testing was corrected using the
Benjamini-Hochberg FDR procedure. Storey's $\pi_1$ was estimated as
$1 - 2 \times P(\text{p-value} > 0.5)$.

% ------------------------------------------------------------
\subsection*{TF motif scanning}
% ------------------------------------------------------------

We used 259 position weight matrices (PWMs) from PlantTFDB~5.0
\cite{tian2020}, representing 36 TF families in \emph{Zea mays}. Motif
scanning used a log-odds score threshold of 6.0 over the $\pm$28\,bp region
centered on each bQTL, requiring physical overlap between the motif match
and the variant nucleotide. This variant-overlap constraint avoids inflating
enrichments for motifs that are common in regulatory regions but do not
intersect the variant \cite{banf2025}.

% ------------------------------------------------------------
\subsection*{TF expression quantification}
% ------------------------------------------------------------

TF genes were identified by matching gene IDs against PlantTFDB~5.0 family
assignments and functional annotations from the EntAP pipeline
\cite{hart2020} applied to B73 NAM5 gene models. Per-family expression was
computed as the mean CPM across all TF genes assigned to each family, using
the well-watered and drought expression datasets separately. Expression
log$_2$FC (drought/WW) was computed per family.

% ------------------------------------------------------------
\subsection*{Condition-dependent switching classification}
% ------------------------------------------------------------

Each bQTL--gene pair was classified as: \emph{constitutive} (FDR $< 0.05$ in
both WW and DS), \emph{WW-specific} (FDR $< 0.05$ in WW only),
\emph{drought-specific} (FDR $< 0.05$ in DS only), or \emph{non-functional}
(not significant in either condition). Gene-level switching was assessed by
aggregating the categories of all bQTL within 2\,kb of each gene's TSS.

% ------------------------------------------------------------
\subsection*{Statistical analysis}
% ------------------------------------------------------------

All statistical tests were two-sided unless otherwise specified. Correlations
were computed as Spearman's $\rho$. Group comparisons used the Mann-Whitney
$U$ test. Enrichments were quantified as odds ratios from $2 \times 2$
contingency tables with Fisher's exact test. Multiple testing correction
used the Benjamini-Hochberg FDR. Analyses were performed in Python~3.14
using scipy, pandas, numpy, and statsmodels.


% ============================================================
\section*{Data and code availability}
% ============================================================

All bQTL, ASE, and expression data are from Engelhorn et al.\
\cite{engelhorn2025} (NCBI GEO: GSE294039, SRA: PRJNA1101486). NAM founder
genotypes are from Grzybowski et al.\ \cite{grzybowski2023} (CyVerse Data
Store). Analysis scripts are available at
\texttt{[repository URL to be added upon publication]}.


% ============================================================
\section*{Acknowledgments}
% ============================================================

We thank Engelhorn et al.\ for making their comprehensive pan-cistrome and
ASE datasets publicly available. This work was supported by [funding sources].


% ============================================================
% REFERENCES
% ============================================================

\begin{thebibliography}{99}

\bibitem{engelhorn2025}
Engelhorn J, Savadel SD, Danz A, et al.
A maize pan-cistrome documents cis-regulatory variation underlying
genotype-specific gene regulation.
\textit{Nature Genetics}. 2025;57:1459--1472.

\bibitem{peleke2024}
Peleke FF, Zumkeller SM, G\"uldener U, Gol L.
DeepCistrome: a deep learning approach for identification of cis-regulatory
elements in plants.
\textit{Nature Communications}. 2024;15:6758.

\bibitem{grzybowski2023}
Grzybowski M, Mural RV, Xu G, et al.
Variation in morpho-physiological traits associated with yield-relevant
source--sink relationships in the NAM founders.
\textit{Crop Science}. 2023;63:3073--3097.

\bibitem{segal2008}
Segal E, Widom J.
From DNA sequence to transcriptional behaviour: a quantitative approach.
\textit{Nature Reviews Genetics}. 2009;10:443--456.

\bibitem{levo2012}
Levo M, Segal E.
In pursuit of design principles of regulatory sequences.
\textit{Nature Reviews Genetics}. 2012;13:786--799.

\bibitem{kumasaka2016}
Kumasaka N, Knights AJ, Gaffney DJ.
Fine-mapping cellular QTLs with RASQUAL and ATAC-seq.
\textit{Nature Genetics}. 2016;48:206--213.

\bibitem{ban2015}
Ban HJ, Heo JY, Oh KS, Park KJ.
Identification of type 2 diabetes-associated combination of SNPs using
support vector machine.
\textit{BMC Genetics}. 2010;11:26.

\bibitem{moyerbrailean2016}
Moyerbrailean GA, Kalita CA, Harvey CT, et al.
Which genetics variants in DNase-I hypersensitive sites are more likely to
alter binding?
\textit{PLoS Genetics}. 2016;12:e1005875.

\bibitem{manolio2009}
Manolio TA, Collins FS, Cox NJ, et al.
Finding the missing heritability of complex diseases.
\textit{Nature}. 2009;461:747--753.

\bibitem{deng2022}
Deng M, Li D, Luo J, et al.
The genetic architecture of amino acids dissection by association and linkage
analysis in maize.
\textit{Plant Biotechnology Journal}. 2017;15:1250--1263.

\bibitem{kim2017}
Kim JM, Sasaki T, Ueda M, Sako K, Seki M.
Chromatin changes in response to drought, salinity, heat, and cold stresses
in plants.
\textit{Frontiers in Plant Science}. 2015;6:114.

\bibitem{tian2020}
Tian F, Yang DC, Meng YQ, Jin J, Gao G.
PlantRegMap: charting functional regulatory maps in plants.
\textit{Nucleic Acids Research}. 2020;48:D1104--D1113.

\bibitem{banf2025}
Banf M, Hartwig T.
An agentic deep learning platform for predicting, explaining, and designing
transcription factor binding variation in maize.
\textit{PLoS Computational Biology}. 2025;(in preparation).

\bibitem{hart2020}
Hart AJ, Ginzburg S, Xu MS, et al.
EnTAP: Bringing faster and smarter functional annotation to non-model
eukaryotic transcriptomes.
\textit{Molecular Ecology Resources}. 2020;20:591--604.

\end{thebibliography}


% ============================================================
% FIGURES
% ============================================================

\clearpage

\begin{figure}[h!]
\centering
% \includegraphics[width=\textwidth]{fig_causal_test.pdf}
\caption{\textbf{Variant genotype causally determines allele-specific expression.}
(\textbf{A})~P-value distribution from genotype--ASE Mann-Whitney tests for 8,040
bQTL--gene pairs under well-watered conditions. The strong left skew (red) relative
to the uniform expectation (dashed line) indicates widespread true associations.
785~pairs are significant at FDR~$< 0.05$ (5.2-fold enrichment).
(\textbf{B})~Effect size distribution for significant pairs (median $|d| = 0.855$).
(\textbf{C})~Example bQTL showing allele-specific expression separated by NAM
founder genotype.}
\label{fig:causal}
\end{figure}

\begin{figure}[h!]
\centering
% \includegraphics[width=\textwidth]{fig_motif_irrelevance.pdf}
\caption{\textbf{TF motif identity is irrelevant under baseline conditions.}
(\textbf{A})~Mean genotype--ASE effect size ($|d|$) by disrupted TF family under
well-watered conditions. All 18~families produce nearly identical effects
(CV = 4.6\%).
(\textbf{B})~Number of disrupted motifs at the variant does not predict effect
magnitude ($r_s = 0.02$).}
\label{fig:motif_irrelevance}
\end{figure}

\begin{figure}[h!]
\centering
% \includegraphics[width=\textwidth]{fig_switching.pdf}
\caption{\textbf{Condition-dependent bQTL switching.}
(\textbf{A})~Venn diagram of functional bQTL--gene pairs under well-watered (WW)
and drought (DS) conditions. Only 31\% are shared (Jaccard = 0.314).
(\textbf{B})~Paired effect sizes for bQTL tested in both conditions, colored by
switching category. WW-specific bQTL (blue) cluster below the diagonal;
drought-specific (red) cluster above.
(\textbf{C})~Effect size ratio (DS/WW) distribution by category, showing the
continuous modulation of bQTL effects.
(\textbf{D})~Target gene drought response (log$_2$FC) by switching category.
Drought-specific bQTL target more drought-responsive genes ($p = 4.1 \times 10^{-4}$).}
\label{fig:switching}
\end{figure}

\begin{figure}[h!]
\centering
% \includegraphics[width=\textwidth]{fig_saturation.pdf}
\caption{\textbf{TF concentration landscape gates bQTL functionality.}
(\textbf{A})~Under well-watered conditions, TF family expression is strongly
anti-correlated with bQTL effect size ($\rho = -0.830$, $p = 0.0008$),
consistent with a saturation model.
(\textbf{B})~Under drought, the correlation collapses ($\rho = -0.067$,
$p = 0.84$) as the TF landscape is reorganized.
(\textbf{C})~TF family expression changes under drought (log$_2$FC) and
corresponding changes in bQTL effect ($\Delta|d|$).
(\textbf{D})~TF family enrichment at drought-specific vs WW-specific bQTL.
Drought-specific bQTL are enriched for motifs of drought-induced TFs
(NAC, TCP, GRAS, bZIP).}
\label{fig:saturation}
\end{figure}


\end{document}
